% ==============================================================================
% Plantilla Institucional de Trabajo Terminal — UPIIZ IPN
% Archivo: config/formato.tex
% ==============================================================================
% Configuración de paquetes, tipografía Arial, geometría, interlineado y formato
% normativo institucional. NO MODIFICAR SALVO INDICACIÓN DOCENTE.
% ==============================================================================

% ------------------------------------------------------------------------------
% 1. IDIOMA Y CARACTERES
% ------------------------------------------------------------------------------
\usepackage[spanish,es-tabla,es-noshorthands]{babel}
\usepackage{csquotes}
% Desactivar interferencias de babel-spanish con porcentajes en tablas
\AtBeginDocument{%
  \renewcommand{\%}{\char`\%}%
}

% ------------------------------------------------------------------------------
% 2. TIPOGRAFÍA INSTITUCIONAL: ARIAL
% ------------------------------------------------------------------------------
% Requisito institucional UPIIZ-IPN:
% - Títulos: Arial 16 pt, negrita
% - Subtítulos: Arial 14 pt, negrita
% - Cuerpo de texto: Arial 12 pt
% - Tablas y Figuras: Arial 10 pt
\usepackage{fontspec}

\IfFontExistsTF{Arial}{%
    \setmainfont{Arial}
}{%
    \IfFontExistsTF{Liberation Sans}{%
        \setmainfont{Liberation Sans}
        \typeout{AVISO: Arial no encontrada en el sistema. Usando Liberation Sans.}
    }{%
        \setmainfont{TeX Gyre Heros}
        \typeout{AVISO: Arial no encontrada. Usando TeX Gyre Heros.}
    }
}

% ------------------------------------------------------------------------------
% 3. MÁRGENES Y GEOMETRÍA DE PÁGINA
% ------------------------------------------------------------------------------
% Requisito institucional:
% - Superior: 2.5 cm
% - Inferior: 2.5 cm
% - Derecho: 2.5 cm
% - Izquierdo: 3.0 cm (encuadernación)
\usepackage{geometry}
\geometry{
    letterpaper,
    top=2.5cm,
    bottom=2.5cm,
    right=2.5cm,
    left=3.0cm,
    headheight=15pt,
    headsep=12pt,
    footskip=20pt
}

% ------------------------------------------------------------------------------
% 4. INTERLINEADO Y JUSTIFICACIÓN
% ------------------------------------------------------------------------------
% Requisito institucional: Interlineado 1.5 y texto justificado.
\usepackage{setspace}
\setstretch{1.5}
\setlength{\parindent}{1.0cm}
\setlength{\parskip}{0pt}

% ------------------------------------------------------------------------------
% 5. FORMATO DE TÍTULOS Y SUBTÍTULOS (titlesec)
% ------------------------------------------------------------------------------
\usepackage{titlesec}

% Títulos de Capítulos: Arial 16 pt, negrita
\titleformat{\chapter}[hang]
    {\fontsize{16}{20}\bfseries\sffamily}
    {\thechapter.\ }
    {0pt}
    {\fontsize{16}{20}\bfseries\sffamily}

\titlespacing*{\chapter}{0pt}{-10pt}{15pt}

% Títulos de Secciones (Subtítulos): Arial 14 pt, negrita
\titleformat{\section}
    {\fontsize{14}{18}\bfseries\sffamily}
    {\thesection\ }
    {0pt}
    {\fontsize{14}{18}\bfseries\sffamily}

\titlespacing*{\section}{0pt}{12pt}{6pt}

% Títulos de Subsecciones: Arial 12 pt, negrita
\titleformat{\subsection}
    {\fontsize{12}{15}\bfseries\sffamily}
    {\thesubsection\ }
    {0pt}
    {\fontsize{12}{15}\bfseries\sffamily}

\titlespacing*{\subsection}{0pt}{10pt}{4pt}

% Títulos de Sub-subsecciones: Arial 12 pt
\titleformat{\subsubsection}
    {\fontsize{12}{15}\itshape\sffamily}
    {\thesubsubsection\ }
    {0pt}
    {\fontsize{12}{15}\itshape\sffamily}

\titlespacing*{\subsubsection}{0pt}{8pt}{3pt}

% ------------------------------------------------------------------------------
% 6. FORMATO DE TABLAS Y FIGURAS
% ------------------------------------------------------------------------------
% Requisito institucional: Arial 10 pt para pies de figura y encabezados de tabla.
\usepackage{caption}
\captionsetup{
    font={small},           % 10 pt
    labelfont={bf,small},   % Etiqueta en negrita
    labelsep=period,        % Punto tras el número: "Figura 1.1."
    format=hang,
    justification=justified
}

\usepackage{subcaption}
\usepackage{booktabs}
\usepackage{tabularx}
\usepackage{longtable}
\usepackage{multirow}
\usepackage{array}
\usepackage{float}
\usepackage{graphicx}
\usepackage{pdfpages}

% ------------------------------------------------------------------------------
% 7. MATEMÁTICAS Y UNIDADES FÍSICAS
% ------------------------------------------------------------------------------
\usepackage{amsmath}
\usepackage{amssymb}
\usepackage{amsfonts}
\usepackage{mathtools}
\usepackage{siunitx}
\sisetup{
    output-decimal-marker = {.},
    per-mode = symbol,
    range-phrase = { a },
    range-units = single
}

% ------------------------------------------------------------------------------
% 8. LISTAS Y ENUMERACIONES
% ------------------------------------------------------------------------------
\usepackage{enumitem}
\setlist{leftmargin=1.5cm, noitemsep, topsep=4pt}

% ------------------------------------------------------------------------------
% 9. ENCABEZADOS Y PIES DE PÁGINA (fancyhdr)
% ------------------------------------------------------------------------------
\usepackage{fancyhdr}
\pagestyle{fancy}
\fancyhf{}
\renewcommand{\headrulewidth}{0pt}
\cfoot{\thepage}

% Estilo para páginas iniciales de capítulo
\fancypagestyle{plain}{%
    \fancyhf{}
    \cfoot{\thepage}
    \renewcommand{\headrulewidth}{0pt}
}

% ------------------------------------------------------------------------------
% 10. HIPERVÍNCULOS Y METADATOS PDF
% ------------------------------------------------------------------------------
\usepackage[
    colorlinks=true,
    linkcolor=black,
    citecolor=black,
    urlcolor=black,
    bookmarks=true,
    bookmarksnumbered=true,
    pdfdisplaydoctitle=true,
    hidelinks
]{hyperref}
\usepackage{bookmark}

% ------------------------------------------------------------------------------
% 11. CÓDIGO FUENTE (listings)
% ------------------------------------------------------------------------------
\usepackage{listings}
\usepackage{xcolor}

\definecolor{codegreen}{rgb}{0,0.6,0}
\definecolor{codegray}{rgb}{0.5,0.5,0.5}
\definecolor{codepurple}{rgb}{0.58,0,0.82}
\definecolor{backcolour}{rgb}{0.97,0.97,0.97}

\lstdefinestyle{estilocodigo}{
    backgroundcolor=\color{backcolour},
    commentstyle=\color{codegreen},
    keywordstyle=\color{blue}\bfseries,
    numberstyle=\tiny\color{codegray},
    stringstyle=\color{codepurple},
    basicstyle=\ttfamily\fontsize{9}{11}\selectfont,
    breakatwhitespace=false,
    breaklines=true,
    captionpos=b,
    keepspaces=true,
    numbers=left,
    numbersep=5pt,
    showspaces=false,
    showstringspaces=false,
    showtabs=false,
    tabsize=2,
    frame=single,
    rulecolor=\color{codegray}
}
\lstset{style=estilocodigo}
